\chapter{Elliptic and hyperbolic geometry}

\begin{convention}[Constant-curvature geometries]
\label{thurston-ch02-constant-curvature}\group{foundations}\source{thurston:page-0018}\notready
The elliptic, Euclidean, and hyperbolic geometries are the complete simply connected Riemannian geometries of constant sectional curvature $+1$, $0$, and $-1$, respectively.
\end{convention}

\begin{definition}[Elliptic space]
\label{thurston-ch02-elliptic-space}\group{elliptic}\source{thurston:page-0018}\notready
Elliptic $n$-space is the quotient of the unit sphere by the antipodal involution, equivalently projective $n$-space with the induced metric. Every two points determine a unique line and the maximum distance is $\pi/2$.
\end{definition}

\begin{definition}[Poincare disk model]
\label{thurston-ch02-poincare-disk}\group{models}\source{thurston:page-0019,thurston:page-0020}\notready
Let $D^n=\{x\in\R^n:\lVert x\rVert<1\}$. In the Poincare disk model of $H^n$, hyperbolic lines and $k$-planes are Euclidean spheres of dimensions $1$ and $k$ orthogonal to $\partial D^n$, with affine planes as limiting cases. The model is conformal and has metric $ds^2=dx^2/(1-r^2)^2$, where $r=\lVert x\rVert$.
\end{definition}

\begin{proposition}[Sphere at infinity]
\label{thurston-ch02-sphere-at-infinity}\group{models}\source{thurston:page-0020}\notready
The boundary $\partial D^n$ is the sphere at infinity. For every base point, geodesic rays determine points of its visual sphere; after moving the base point to the disk origin, this visual sphere is identified with $\partial D^n$ by Euclidean ray endpoints.
\end{proposition}
\begin{proof}\sketch A ray and its base point span a hyperbolic two-plane, hence determine a great-circle arc in the visual sphere. Its endpoint gives the corresponding boundary point, and translation to the origin identifies it with the disk endpoint.\end{proof}

\begin{definition}[Southern-hemisphere model]
\label{thurston-ch02-southern-hemisphere}\group{models}\source{thurston:page-0020,thurston:page-0021}\notready
The inclusion $D^n\subset D^{n+1}$ is a hyperbolic $n$-plane in $H^{n+1}$. Euclidean stereographic projection from the north pole sends it to the southern hemisphere; hyperbolic lines become circles orthogonal to the equator. Intrinsically, the image of $p$ is the endpoint at infinity of the ray perpendicular to the hyperplane at $p$ and directed downward.
\end{definition}

\begin{definition}[Upper half-space model]
\label{thurston-ch02-upper-half-space}\group{models}\source{thurston:page-0021,thurston:page-0022}\notready
The upper half-space model is $\{x_n>0\}\subset\R^n$ with metric $ds^2=dx^2/x_n^2$. Its hyperbolic lines are circles orthogonal to the boundary plane $\R^{n-1}$, including vertical lines; Euclidean size is proportional to height above that boundary.
\end{definition}

\begin{definition}[Projective model]
\label{thurston-ch02-projective-model}\group{models}\source{thurston:page-0022,thurston:page-0023}\notready
Orthogonal projection of the southern hemisphere to its equatorial disk gives the projective model. Hyperbolic lines are Euclidean line segments, so incidence is preserved but angles and shapes are not. Two lines are parallel when they meet on the sphere at infinity; otherwise they are ultraparallel, with a unique common perpendicular. Their projective intersection lies outside the sphere at infinity, and every line through it is perpendicular to that common perpendicular.
\end{definition}
\begin{proof}\sketch Put $H^2$ in a plane $P\subset H^3$ and let $Q$ be the plane through the common perpendicular $L$ perpendicular to $P$. The plane through an observer and any line of $P$ perpendicular to $L$ is perpendicular to $Q$, so its visual line passes through the visual point of $L$; the converse gives the claim.\end{proof}

\begin{proposition}[Projective duality for hyperplanes]
\label{thurston-ch02-projective-duality}\group{models}\source{thurston:page-0025}\notready
Points outside the sphere at infinity correspond to hyperplanes of $H^n$ via the common intersection point of all their perpendiculars. Dually, a point of $H^n$ corresponds to the locus of points outside the sphere at infinity represented by hyperplanes through it.
\end{proposition}

\begin{definition}[Hyperboloid model]
\label{thurston-ch02-hyperboloid}\group{models}\source{thurston:page-0025,thurston:page-0026,thurston:page-0027}\notready
On $\R^{n+1}$ use $Q(X)=X_1^2+\cdots+X_n^2-X_{n+1}^2$. The positive sheet $H^n=\{Q=-1,X_{n+1}>0\}$, with the restricted form, is hyperbolic space of curvature $-1$. Planes through the origin meet it in geodesics. Projection from the origin gives the projective model, and projective transformations preserving the sphere at infinity are induced by isometries of $Q$; hence hyperbolic space has no non-isometric similarities.
\end{definition}

\begin{notation}[Lorentzian cosine]
\label{thurston-ch02-lorentzian-cosine}\group{trigonometry}\source{thurston:page-0027}\notready
Write $X\cdot Y=\sum_{i=1}^nX_iY_i-X_{n+1}Y_{n+1}$ and $c(X,Y)=(X\cdot Y)/(\sqrt{X\cdot X}\sqrt{Y\cdot Y})$ for non-null vectors. A real-length vector $Y$ represents the oriented hyperplane $Y^\perp$.
\end{notation}

\begin{proposition}[Geodesic equation on a level set]
\label{thurston-ch02-geodesic-equation}\group{trigonometry}\source{thurston:page-0027}\dcref{2.6.1}\notready
If $X_t$ is a geodesic on $Q=r^2$ with speed $s=\sqrt{Q(\dot X_t)}$, then $\ddot X_t=-(s/r)^2X_t$, with $X_0\cdot X_0=r^2$, $\dot X_0\cdot\dot X_0=s^2$, and $X_0\cdot\dot X_0=0$.
\end{proposition}
\begin{proof}\sketch Differentiate the constant-norm equations to obtain the three orthogonality identities; since a geodesic lies in the span of its initial position and velocity, they force the displayed acceleration.\end{proof}

\begin{theorem}[Distance from the Lorentzian cosine]
\label{thurston-ch02-cosh-distance}\group{trigonometry}\source{thurston:page-0028}\dcref{2.6.2}\notready
For $X,Y\in H^n$, $c(X,Y)=\cosh d(X,Y)$.
\end{theorem}
\begin{proof}\sketch Apply the geodesic equation to the unit-speed geodesic from $X$ to $Y$; its inner product with $X$ is $\cosh t$, and evaluation at $t=d(X,Y)$ proves the formula.\end{proof}

\begin{proposition}[Incidence and angle classification]
\label{thurston-ch02-hyperplane-classification}\group{trigonometry}\source{thurston:page-0028,thurston:page-0029}\dcref{2.6.3,2.6.4,2.6.5}\notready
For distinct oriented hyperplanes $X^\perp,Y^\perp$, they intersect iff $Q$ is positive definite on $\langle X,Y\rangle$, equivalently $c(X,Y)^2<1$, and then $c(X,Y)=\cos\angle(X,Y)$. They have a common perpendicular iff the restriction has type $(1,1)$, equivalently $c(X,Y)^2>1$, and then $c(X,Y)=\pm\cosh d(X^\perp,Y^\perp)$, with sign determined by normal orientations. They are parallel iff the restriction is degenerate, equivalently $c(X,Y)^2=1$, and their set-distance is zero.
\end{proposition}
\begin{proof}\sketch The signature on the normal span is complementary to that on $X^\perp\cap Y^\perp$. Its Gram determinant is $\lVert X\rVert^2\lVert Y\rVert^2(1-c(X,Y)^2)$, giving the three cases; the geodesic equation in the $(1,1)$ case gives the common-perpendicular distance relation, and the degenerate case is its limit.\end{proof}

\begin{proposition}[Point-hyperplane distance]
\label{thurston-ch02-point-hyperplane}\group{trigonometry}\source{thurston:page-0030}\dcref{2.6.6}\notready
If $X\in H^n$ and $Y^\perp$ is oriented, then $c(X,Y)=\sinh(\dist(X,Y^\perp))/i$, where the distance is oriented.
\end{proposition}

\begin{proposition}[Dual-basis cosine formula]
\label{thurston-ch02-dual-cosine}\group{trigonometry}\source{thurston:page-0030,thurston:page-0031}\dcref{2.6.7}\notready
For projective triangle representatives $v_1,v_2,v_3$, put $\epsilon_i=v_i\cdot v_i$, $\epsilon_{ij}=\sqrt{\epsilon_i\epsilon_j}$, and $c_{ij}=c(v_i,v_j)$. If $C_{ii}=\epsilon_i$, $C_{ij}=\epsilon_{ij}c_{ij}$, and $v^1,v^2,v^3$ is the dual basis, then
\[
c^{12}=\epsilon\frac{c_{13}c_{23}-c_{12}}{\sqrt{1-c_{23}^2}\sqrt{1-c_{13}^2}},\qquad
\epsilon=\frac{-\epsilon_{12}\epsilon_3}{\sqrt{-\epsilon_2\epsilon_3}\sqrt{-\epsilon_1\epsilon_3}},
\]
where $c^{ij}=c(v^i,v^j)$.
\end{proposition}
\begin{proof}\sketch The dual Gram matrix is $C^{-1}$. Compute it from $-\operatorname{adj}C=(-\det C)C^{-1}$, normalize the $(1,2)$ entry, and determine its sign from the diagonal norms.\end{proof}

\begin{proposition}[Hyperbolic triangle identities]
\label{thurston-ch02-triangle-identities}\group{trigonometry}\source{thurston:page-0031,thurston:page-0033,thurston:page-0034}\dcref{2.6.8,2.6.9,2.6.11,2.6.12,2.6.13,2.6.14,2.6.15,2.6.16}\notready
For a triangle with sides $A,B,C$ opposite angles $\alpha,\beta,\gamma$,
\[
\cosh C=\frac{\cos\alpha\cos\beta+\cos\gamma}{\sin\alpha\sin\beta},\quad
\cos\gamma=\frac{\cosh A\cosh B-\cosh C}{\sinh A\sinh B},\quad
\cosh C=\cosh A\cosh B-\sinh A\sinh B\cos\gamma,
\]
and $\sinh A/\sin\alpha=\sinh B/\sin\beta=\sinh C/\sin\gamma$. If the triangle is right at $\gamma$, then $\cosh C=\cosh A\cosh B$, $\cosh A=\cos\alpha/\sin\beta$, and $\sin\alpha=\sinh A/\sinh C$. At an ideal vertex the first identity becomes $\cosh C=(\cos\alpha\cos\beta+1)/(\sin\alpha\sin\beta)$, and with a right angle $\cosh C=1/\sin\alpha$.
\end{proposition}
\begin{proof}\sketch Represent vertices by Lorentzian vectors, apply the dual-basis formula, and specialize the Gram entries. Taking a vertex to the sphere at infinity gives the limiting identities; setting an angle to $\pi/2$ gives the right-triangle formulas.\end{proof}

\begin{proposition}[Right polygon identities]
\label{thurston-ch02-right-polygon-identities}\group{trigonometry}\source{thurston:page-0031,thurston:page-0032,thurston:page-0035}\dcref{2.6.10,2.6.17,2.6.18}\notready
For an all-right hexagon,
\[
\cosh C=\frac{\cosh\alpha\cosh\beta+\cosh\gamma}{\sinh\alpha\sinh\beta},\qquad
\frac{\sinh A}{\sinh\alpha}=\frac{\sinh B}{\sinh\beta}=\frac{\sinh C}{\sinh\gamma}.
\]
For an all-right pentagon, $\sinh A\sinh B=\cosh D$.
\end{proposition}
\begin{proof}\sketch Apply the dual-basis formula to decompositions into right triangles and eliminate intermediate diagonals; the pentagon identity is the corresponding two-triangle calculation, and comparison of the three hexagon expressions gives its sine law.\end{proof}
